\section{Dirichlet's theorem with an explicit primitive recursive bound}
\label{sec:direct-dirichlet}

We give a quantitative proof of Dirichlet's theorem which does not use a
prime number theorem, a zero-free region, a class number formula, or an
ineffective lower bound for an $L$-value.  Its central step is a finite
coefficient argument for real characters.  In particular, the argument never
evaluates a Dirichlet series outside its half-plane of convergence.  All
uses of infinite sums below come with bounds that permit replacement by
finite rational calculations, as specified in the arithmetic implementation
section.

\begin{theorem}[An explicit bound for Dirichlet's theorem]
\label{thm:dirichlet-explicit}
Let $q\geq1$, $a,B\in\mathbb N$, and $\gcd(a,q)=1$. Define the integers
\begin{align*}
 D&=2^{25}(q+1)^2,\\
 C&=1+16q^2 2^D,\\
 T&=8q(B+C+q+4),\\
 Y&=2^{T2^T}.
\end{align*}
Then
\[
 \sum_{\substack{p\leq Y\\p\text{ prime}\\p\equiv a\pmod q}}
       \frac1p>B+1.
\]
Consequently there is a prime $p$ with
\[
 B<p\leq Y,\qquad p\equiv a\pmod q.
\]
The bound $Y(q,B)$ is primitive recursive.  Bounded search through the
integers up to $Y(q,B)$ therefore gives a primitive recursive next-prime
function for every reduced residue class, with no appeal to an unbounded
search procedure.
\end{theorem}

The constants are intentionally large.  Their role is to make every
existential choice quantitatively explicit.  The case $q=1$ is included:
the character group then has one element, and assertions about
nonprincipal characters are empty assertions.

\subsection{Finite characters and orthogonality}

Let $G=(\mathbb Z/q\mathbb Z)^\times$ and let $h=|G|=\varphi(q)$.
A character is a homomorphism from $G$ to the complex roots of unity.
Since every element of $G$ has order dividing $h$, all character values are
$h$-th roots of unity.  Thus characters can be represented by finite tables:
enumerate the at most $h^h$ maps from $G$ to the $h$-th roots of unity and
retain the maps satisfying the multiplication-table equations.  This is
a bounded construction.

Here is a finite proof that there are exactly $h$ characters and that they
separate the elements of $G$.  Suppose that a character $\chi_H$ is defined
on a subgroup $H\subseteq G$, and choose $g\notin H$.  Let $m$ be the least
positive integer with $g^m\in H$.  The subgroup $H'=\langle H,g\rangle$
has the unique representatives $ug^j$ with $u\in H$ and $0\leq j<m$.
Choose a root $\alpha$ of
\[
 \alpha^m=\chi_H(g^m).
\]
There are exactly $m$ choices.  Because $\chi_H(g^m)$ is an $|H|$-th root
of unity, each such $\alpha$ is an $m|H|=|H'|$-th root of unity, hence an
$h$-th root of unity.  Define
\[
 \chi_{H'}(ug^j)=\chi_H(u)\alpha^j.
\]
Uniqueness of the representatives and the equation for $\alpha$ prove that
this is a character.  Conversely every extension has this form.  Beginning
with the trivial subgroup and adjoining elements until $G$ is reached
therefore gives exactly $|G|$ characters.  To separate $v\neq1$, first
give the cyclic subgroup generated by $v$ the character taking $v$ to a
primitive root of order $\operatorname{ord}(v)$, and then extend it as above.

If $v\neq1$, choose a character $\chi_*$ with $\chi_*(v)\neq1$.
Multiplication by $\chi_*$ permutes the character list, so
\[
 \sum_\chi\chi(v)=\chi_*(v)\sum_\chi\chi(v),
 \qquad\text{hence}\qquad \sum_\chi\chi(v)=0.
\]
At $v=1$ the sum is $h$.  Extending a character by zero on the integers not
coprime to $q$ gives the usual Dirichlet characters, with the orthogonality
identity
\begin{equation}
 \sum_\chi\overline{\chi(a)}\chi(n)
   =h\,\mathbf 1_{n\equiv a\, (\mathrm{mod}\ q)}.
 \label{eq:finite-character-orthogonality}
\end{equation}
For a nonprincipal character, summing over one period gives zero: multiply
the summands by a residue on which the character is not $1$ and permute the
reduced residues.  Consequently its partial sums satisfy
\begin{equation}
 A_\chi(t):=\sum_{1\leq n\leq t}\chi(n),\qquad |A_\chi(t)|\leq q.
 \label{eq:partial-character-bound}
\end{equation}
The finite root choices and the subgroup extensions in this paragraph can
all be obtained from their bounded lists.  No choice principle for an
infinite family is involved.

\subsection{Elementary analytic estimates for a nonprincipal character}

For a nonprincipal character, partial summation gives
\begin{equation}
 L(s,\chi)=s\int_1^\infty A_\chi(t)t^{-s-1}\,dt\qquad(s\geq1).
 \label{eq:L-partial-integral}
\end{equation}
The integral is an abbreviation for effectively convergent integrals over
finite intervals.  It also defines the endpoint $L(1,\chi)$; for example,
the tail of the corresponding Dirichlet series after $N$ is at most
$2q/N$.  Equation~\eqref{eq:partial-character-bound} gives
\begin{equation}
 |L(s,\chi)|\leq q,\qquad |L'(s,\chi)|\leq\frac{2q}{s}\leq2q
 \quad(s\geq1).
 \label{eq:L-basic-estimates}
\end{equation}
Indeed the second inequality follows by differentiating
\eqref{eq:L-partial-integral} and bounding the two terms by
\[
 q\int_1^\infty t^{-s-1}(1+s\log t)\,dt=\frac{2q}{s}.
\]
Differentiation is justified by the same explicitly integrable majorants.
In particular,
\begin{equation}
 |L(1+t,\chi)-L(1,\chi)|\leq2qt\qquad(t\geq0).
 \label{eq:L-lipschitz}
\end{equation}

We will also use, for $s>1$, the elementary estimates
\begin{equation}
 \frac1{s-1}\leq\zeta(s)\leq1+\frac1{s-1}.
 \label{eq:zeta-integral-test}
\end{equation}
They are the lower and upper integral comparisons for the decreasing
function $t^{-s}$.  Finite unique factorization proves the Euler products
for $\zeta$ and $L$; their tails are dominated by the same convergent
positive series.  Define the Euler logarithm by the particular convergent
series
\begin{equation}
 F_\chi(s)=\sum_p\sum_{j\geq1}\frac{\chi(p)^j}{j p^{js}},
 \qquad s>1.
 \label{eq:Euler-log}
\end{equation}
Then $\exp F_\chi(s)=L(s,\chi)$ and
$F'_\chi(s)=L'(s,\chi)/L(s,\chi)$.  Defining this logarithm by its series
avoids any choice of a complex logarithm branch.

\subsection{The real-character nonvanishing estimate}

\begin{lemma}
\label{lem:real-L-margin}
For every real nonprincipal character modulo $q$,
\[
 L(1,\chi)>2^{-D},\qquad D=2^{25}(q+1)^2.
\]
\end{lemma}

\begin{proof}
All character values are now in $\{-1,0,1\}$.  Write $\varepsilon=L(1,\chi)$.
We will obtain a positive lower bound directly, without assuming
$\varepsilon\neq0$.

First expand the integral~\eqref{eq:L-partial-integral} around $s=2$:
\[
 L(2-z,\chi)=\sum_{k\geq0}l_kz^k.
\]
The coefficient formulas are
\begin{align*}
 l_0&=2\int_1^\infty A_\chi(t)t^{-3}\,dt,\\
 l_k&=\int_1^\infty A_\chi(t)t^{-3}
 \left(\frac{2(\log t)^k}{k!}
       -\frac{(\log t)^{k-1}}{(k-1)!}\right)\,dt\quad(k\geq1).
\end{align*}
Since
\[
 \int_1^\infty t^{-3}(\log t)^k\,dt=\frac{k!}{2^{k+1}},
\]
we have
\begin{equation}
 |l_0|\leq q,\qquad |l_k|\leq\frac{2q}{2^k}\quad(k\geq1).
 \label{eq:l-coefficients}
\end{equation}
The exponential-series expansion inside the integral is dominated on every
$|z|\leq r_0<2$, and these coefficient estimates also give its explicit
geometric tails.  In particular $\varepsilon=\sum_{k\geq0}l_k$.

For $s>1$ the elementary summation formula for $\zeta$ gives
\[
 Z(s):=\zeta(s)-\frac1{s-1}
       =1-s\int_1^\infty\{t\}t^{-s-1}\,dt.
\]
Expand the right-hand side as
$Z(2-z)=\sum_{k\geq0}u_kz^k$.  Using $0\leq\{t\}<1$ in exactly the preceding
calculation gives
\begin{equation}
 |u_0|\leq2,\qquad |u_k|\leq\frac2{2^k}\quad(k\geq1).
 \label{eq:u-coefficients}
\end{equation}
These integral coefficient formulas require only real-variable integration
and the exponential series.

For $|z|<1$ write
\begin{align}
 \zeta(2-z)L(2-z,\chi)
   &=\frac{\varepsilon}{1-z}+H(z),\label{eq:positive-series-decomp}\\
 H(z)&=Z(2-z)L(2-z,\chi)
       +\frac{L(2-z,\chi)-\varepsilon}{1-z}.
 \nonumber
\end{align}
If the coefficients of the second term in $H$ are $b_k$, then
\[
 b_k=\sum_{j=0}^k l_j-\varepsilon=-\sum_{j>k}l_j,
 \qquad |b_k|\leq\frac{2q}{2^k}.
\]
Consequently the coefficients $h_k$ of $H$ satisfy
\begin{equation}
 |h_k|\leq\frac{q(4k+6)}{2^k}\qquad(k\geq0).
 \label{eq:h-coefficients}
\end{equation}
For example, bound each of the $k+1$ terms of the product convolution by
$4q/2^k$ and add the bound for $b_k$.

On the other hand, the absolutely convergent Dirichlet product at $s>1$ is
\[
 \zeta(s)L(s,\chi)=\sum_{n\geq1}\frac{a_n}{n^s},
 \qquad a_n=\sum_{d\mid n}\chi(d).
\]
The coefficients $a_n$ are multiplicative.  For a prime power $p^e$ the
coefficient is $e+1$ when $\chi(p)=1$, is $1$ when $\chi(p)=0$, and is
$1$ or $0$ according as $e$ is even or odd when $\chi(p)=-1$.
It follows that
\begin{equation}
 a_n\geq0\quad\text{for all }n,\qquad a_{m^2}\geq1\quad\text{for all }m.
 \label{eq:positive-square-coefficients}
\end{equation}
Its coefficients about $s=2$ are therefore
\begin{equation}
 c_k=\frac1{k!}\sum_{n\geq1}\frac{a_n(\log n)^k}{n^2}\geq0.
 \label{eq:c-coefficients}
\end{equation}
Expanding on any sufficiently small fixed interval around $z=0$ and
comparing coefficients in~\eqref{eq:positive-series-decomp} gives
\begin{equation}
 c_k=\varepsilon+h_k\qquad(k\geq0).
 \label{eq:c-epsilon-h}
\end{equation}
This coefficient comparison can equivalently be obtained by differentiating
the convergent series at $s=2$.  Its convergence is quantitative: since
$a_n\leq d(n)\leq2\sqrt n$ and
$(\log n)^k/k!\leq4^k n^{1/4}$, the tail after $N$ in
\eqref{eq:c-coefficients} is at most
\begin{equation}
 8\,4^kN^{-1/4}.
 \label{eq:c-coefficient-tail}
\end{equation}

Now form a \emph{finite} polynomial.  Put
\[
 r=\frac74,\qquad M=2^{20}(q+1)^2,\qquad K=16M,
 \qquad P_K=\sum_{k=0}^K c_kr^k.
\]
From~\eqref{eq:h-coefficients} and~\eqref{eq:c-epsilon-h},
\begin{align}
 P_K
 &\leq\varepsilon\sum_{k=0}^K r^k
       +q\sum_{k\geq0}(4k+6)(7/8)^k\nonumber\\
 &=\varepsilon\sum_{k=0}^K r^k+272q.
 \label{eq:PK-upper}
\end{align}
Only the absolutely convergent majorant for $h_k$ was summed at $r$.
We have made no assertion that the Dirichlet series for $\zeta L$ converges
there.

For $t\geq0$ let $E_K(t)=\sum_{k=0}^Kt^k/k!$.  Its exponential tail obeys
\[
 e^t-E_K(t)
 \leq\frac1{K+1}\sum_{j>K}\frac{jt^j}{j!}
 \leq\frac{t}{K+1}e^t.
\]
In particular $K\geq4t$ implies $E_K(t)\geq\tfrac12e^t$.
For $m\leq M$ set $t_m=2r\log m=\tfrac72\log m$.  Since
$t_m\leq\tfrac72M$ and $K=16M$, this estimate applies.  Retain only the
square terms of~\eqref{eq:c-coefficients}, using
\eqref{eq:positive-square-coefficients}, to obtain
\begin{align}
 P_K
 &\geq\sum_{m=1}^M\frac1{m^4}
                \sum_{k=0}^K\frac{(2r\log m)^k}{k!}\nonumber\\
 &\geq\frac12\sum_{m=1}^M m^{-1/2}
 \geq\frac12\sqrt M=512(q+1).
 \label{eq:PK-lower}
\end{align}
Combining~\eqref{eq:PK-upper} and~\eqref{eq:PK-lower} gives
\[
 \varepsilon\sum_{k=0}^K r^k\geq240q+512>1.
\]
As $r<2$, the finite geometric sum is smaller than $2^{K+1}$.  Hence
$\varepsilon>2^{-K-1}$.  Finally
$D=2^{25}(q+1)^2=2K\geq K+1$, proving the claimed lower bound.
\end{proof}

The finite polynomial $P_K$ is the reason this proof has a usable
constructive form.  Positivity supplies a large finite lower bound from the
squares, while the coefficient estimates supply a finite upper bound.  No
principle asserting the existence of a singularity at an unspecified
abscissa is needed.

\subsection{The nonreal-character estimate}

\begin{lemma}
\label{lem:nonreal-L-margin}
For a nonreal character modulo $q$,
\[
 |L(1,\chi)|\geq\frac1{2048q^4}\geq2^{-D}.
\]
\end{lemma}

\begin{proof}
The character $\chi^2$ is nonprincipal: otherwise all nonzero values of
$\chi$ would be $1$ or $-1$.  For a complex number $z$ of modulus one,
\[
 3+4\operatorname{Re}z+\operatorname{Re}(z^2)
   =2(1+\operatorname{Re}z)^2\geq0.
\]
Apply this inequality term by term to the Euler logarithms.  At primes
dividing $q$ the corresponding numerator is $3$, which is also nonnegative.
Exponentiation gives, for $s>1$,
\begin{equation}
 \zeta(s)^3|L(s,\chi)|^4|L(s,\chi^2)|\geq1.
 \label{eq:nonreal-product-positive}
\end{equation}
Set $t=1/(4096q^5)$ and $s=1+t$.  We have
$\zeta(s)\leq2/t$ by~\eqref{eq:zeta-integral-test}, and
$|L(s,\chi^2)|\leq q$ by~\eqref{eq:L-basic-estimates}.  Thus
\[
 |L(1+t,\chi)|^4\geq\frac{t^3}{8q}
     =\frac1{2^{39}q^{16}}\geq\frac1{(1024q^4)^4}.
\]
Using~\eqref{eq:L-lipschitz},
\[
 |L(1,\chi)|
 \geq\frac1{1024q^4}-2qt
 =\frac1{2048q^4}.
\]
Finally $q\leq2^q$ and $D\geq11+4q$ imply
$2048q^4\leq2^D$.
\end{proof}

\subsection{A uniform bound for the Euler logarithms}

Put $\delta=2^{-D}$.  The two preceding lemmas give
\begin{equation}
 |L(1,\chi)|\geq\delta
 \quad\text{for every nonprincipal character }\chi.
 \label{eq:uniform-L-margin}
\end{equation}
We next claim that
\begin{equation}
 |L(s,\chi)|\geq\frac\delta{8q}
 \qquad(1\leq s\leq2).
 \label{eq:L-interval-margin}
\end{equation}
Write $s=1+t$.  If $0\leq t\leq\delta/(4q)$,
\eqref{eq:L-lipschitz} and~\eqref{eq:uniform-L-margin} give
$|L(s,\chi)|\geq\delta/2$.  If $\delta/(4q)\leq t\leq1$, the Euler product
gives $|L(s,\chi)|\geq1/\zeta(s)\geq t/2\geq\delta/(8q)$.
For the first inequality, bound each reciprocal Euler factor by
$|1-\chi(p)p^{-s}|\leq1+p^{-s}\leq(1-p^{-s})^{-1}$ and pass to the
effectively convergent products.

At $s=2$, equation~\eqref{eq:Euler-log} gives
$|F_\chi(2)|\leq\log\zeta(2)\leq\log2<1$.  On $1<s\leq2$,
equations~\eqref{eq:L-basic-estimates} and~\eqref{eq:L-interval-margin} give
\[
 |F'_\chi(s)|=\left|\frac{L'(s,\chi)}{L(s,\chi)}\right|
       \leq\frac{16q^2}{\delta}.
\]
Integration over an interval of length at most one yields the explicit bound
\begin{equation}
 |F_\chi(s)|\leq1+16q^2 2^D=C
 \qquad(1<s\leq2)
 \label{eq:Euler-log-bound}
\end{equation}
for every nonprincipal character.

\subsection{Orthogonality produces a large prime sum}

Let
\[
 S_a(s)=\sum_{\substack{p\text{ prime}\\p\equiv a\pmod q}}p^{-s}
 \qquad(s>1).
\]
Apply~\eqref{eq:finite-character-orthogonality} inside the absolutely
convergent Euler logarithms.  The first prime-power terms give $hS_a(s)$;
the remaining terms give a real number $hR_a(s)$ with
\begin{align*}
 0\leq R_a(s)
 &\leq\sum_p\sum_{j\geq2}\frac1{j p^{js}}\\
 &\leq\sum_{n\geq2}\sum_{j\geq2}n^{-j}
   =\sum_{n\geq2}\frac1{n(n-1)}=1.
\end{align*}
Therefore
\begin{equation}
 hS_a(s)=\sum_\chi\overline{\chi(a)}F_\chi(s)-hR_a(s).
 \label{eq:prime-log-orthogonality}
\end{equation}

For the principal character $\chi_0$,
\[
 L(s,\chi_0)=\zeta(s)\prod_{p\mid q}(1-p^{-s})
  \geq\frac1{s-1}\prod_{p\mid q}(1-p^{-1})
  =\frac{\varphi(q)}{q(s-1)}\geq\frac1{q(s-1)}.
\]
Its Euler logarithm is real and so satisfies
\[
 F_{\chi_0}(s)\geq\log\frac1{s-1}-\log q.
\]
Take real parts in~\eqref{eq:prime-log-orthogonality} and use
\eqref{eq:Euler-log-bound}.  The result, in a slightly weakened convenient
form, is
\begin{equation}
 hS_a(s)\geq\log\frac1{s-1}-\log q-h(C+1),
 \qquad 1<s\leq2.
 \label{eq:prime-sum-lower}
\end{equation}

Now choose $T$ and $Y$ as in Theorem~\ref{thm:dirichlet-explicit} and set
$\tau=2^{-T}$, $s=1+\tau$.  Since $\log2\geq1/2$, $h\leq q$, and
$\log q\leq q$, equation~\eqref{eq:prime-sum-lower} gives
\begin{align*}
 hS_a(1+\tau)
 &\geq4q(B+C+q+4)-q-q(C+1)\\
 &=4qB+3qC+4q^2+14q.
\end{align*}
In particular $S_a(1+\tau)\geq B+3$.  The prime tail after $Y$ is bounded
by the integer tail, and the integral test gives
\[
 \sum_{p>Y}p^{-1-\tau}
 \leq\sum_{n>Y}n^{-1-\tau}
 \leq\frac{Y^{-\tau}}\tau=1,
\]
because $Y=2^{T2^T}$.  Hence
\[
 \sum_{\substack{p\leq Y\\p\equiv a\pmod q}}\frac1p
 \geq\sum_{\substack{p\leq Y\\p\equiv a\pmod q}}p^{-1-\tau}
 \geq B+2>B+1.
\]
If all the primes in this finite sum were at most $B$, there would be at
most $B$ summands and each would be at most $1$, contradicting the displayed
inequality.  This completes the proof of
Theorem~\ref{thm:dirichlet-explicit}.

\subsection{What the quantitative proof supplies to the arithmetic layer}

The conclusion is a finite rational inequality and a bounded prime-search
specification.  The auxiliary real and complex quantities serve as readable
abbreviations for approximations.  A few particularly useful schedules are
as follows.

For the real-character coefficient polynomial, equation
\eqref{eq:c-coefficient-tail} gives the simultaneous weighted tail bound
\[
 \sum_{k=0}^K r^k\frac1{k!}
     \sum_{n>N}\frac{a_n(\log n)^k}{n^2}
 \leq2\,7^{K+1}N^{-1/4}.
\]
Thus the entirely integral choice
\[
 N=2^{4(p+4)}7^{4(K+1)}
\]
leaves a tail at most $2^{-p-3}$.  The bounds for $l_k,u_k,b_k,h_k$ have
explicit geometric tails.  The square lower bound in
\eqref{eq:PK-lower} uses only $M$ terms and $K+1$ terms of an exponential
polynomial.  Its gap from the upper bound is much larger than $1$, so fixed
rational approximation allowances can be reserved without changing any
conclusion.

At the final argument $s=1+2^{-T}$, the cutoff
\[
 N_p=2^{(p+T+4)2^T}
\]
gives
$\sum_{n>N_p}n^{-s}\leq2^{-p-4}$.  All these integers are primitive
recursive in the displayed parameters.  For integrations of $L'/L$, the
same partial-summation representation gives $|L''(s)|\leq4q$ on $[1,2]$.
Together with~\eqref{eq:L-interval-margin} this gives the explicit
Lipschitz bound
\[
 \left|\left(\frac{L'}{L}\right)'\right|
 \leq\frac{32q^2}{\delta}+\frac{256q^4}{\delta^2}
 \leq\frac{288q^4}{\delta^2}.
\]
It therefore suffices to use a uniform rational mesh whose size is chosen
from this bound and the requested accuracy.

These schedules identify the finite work needed for each approximation.
To finish the formal iPRA development one proves the associated rational
inequalities by quantifier-free induction, using the coding and
approximation lemmas described elsewhere in this report.  An ordinary
Lean theorem about $\mathbb R$ or $\mathbb C$ is a useful semantic check, but
does not by itself replace those object-language derivations.
